% sections/prevalence.tex -- Section 4: prevalence of under-binding (E1, E2). A1 moved to sections/ablation.tex (W25).
% Owner: W10 (wave 2). Every number is a macro from numbers.tex, whose comments name
% the source line in binding/prevalence/E1_results.md or E2_results.md (W8).
\section{Prevalence of Under-Binding (\GA)}\label{sec:prevalence}

\GA needs no bug: a proof bound to $(\htm,\htu)$ leaves the operation open
whenever the executor dispatches on something outside $\Cov$. Method
confusion and dispatch are known: Golinelli et al.\ show that applications
treat POST and GET interchangeably~\cite{p2:golinelli2026post}, and
CVE-2026-37236 shows an \hdr{X-HTTP-Method-Override} bypassing upstream method
ACLs~\cite{grpcgw2026cve37236}. Our contribution is to measure how often a
proof-bound request is exposed: how often API descriptions put several
operations behind one (method, path) key (\Exp{1}), and how often method
override---which survives a proof-bound \htm because it is resolved after
verification---is reachable (\Exp{2}).

\subsection{API descriptions (\Exp{1})}\label{sec:prev:e1}

\paragraph{Method.} The verifier's view is the key $K = (\textit{method},
\textit{path})$. A key is \emph{collided} if the description maps it to two or
more operations selected by a query parameter, a header or a body field. We
analyse the APIs.guru OpenAPI directory (one pinned commit, cloned 2026-11-22):
\NumSpecFiles{} spec files, of which \NumSpecFilesParsed{} parsed
(\NumSpecFilesFailed{} YAML errors excluded). The unit is one preferred spec per
API, giving \NumSpecs{} APIs with \NumOps{} operations, \NumStateChangingOps{}
of them state-changing (POST, PUT, PATCH or DELETE). Five detectors flag
collisions: a query or fragment inside the path key (M1), a dispatcher-named
parameter with an enum of at least two values (M2), header dispatch such as
\hdr{X-Amz-Target} (M3), a discriminated body union or dispatcher-named body
field (M4), and JSON-RPC or GraphQL single endpoints (M5). Only their strong
tiers are counted.

\paragraph{Validation.} One author judged a seeded, detector-stratified sample of
\NumValidationSample{} flagged operations against the spec text (unclear counts as false positive). Precision is
\PrecMOne{} for M1 (\NumMOneTP/\NumMOneSample), \PrecMTwo{} for M2, \PrecMFour{} for M4 and
\PrecMFive{} for M5 (\NumMFiveFlags{} flags). Weighted by stratum it is \PrecWeighted{}
overall, but only \PrecNonMOneWeighted{} outside M1. We therefore also report a
hand-confirmed lower bound.

\begin{table}[t]
\caption{\Exp{1}: state-changing operations on a collided (method, path) key.
Raw detector counts, strong tier; 95\% Wilson intervals in brackets.}
\label{tab:e1}
\small
\begin{tabular}{@{}l r r r@{}}
\toprule
Subset & APIs & APIs collided & State ops collided \\
\midrule
All & \NumSpecs & \NumApisCollided{} (\PctApisCollided) & \NumStateOpsCollided{} (\PctUnderBound) \\
Excl.\ AWS & \NumSpecsNonAWS & \NumApisCollidedNonAWS{} (\PctApisCollidedNonAWS) & \NumStateOpsCollidedNonAWS{} (\PctUnderBoundNonAWS) \\
Excl.\ hyperscalers & \NumSpecsNonHyper & \NumApisCollidedNonHyper{} (\PctApisCollidedNonHyper) & \NumStateOpsCollidedNonHyper{} (\PctUnderBoundNonHyper) \\
Finance & \NumSpecsFinance & \NumApisCollidedFinance{} (\PctApisCollidedFinance) & \NumStateOpsCollidedFinance{} (\PctUnderBoundFinance) \\
\bottomrule
\end{tabular}
\end{table}

\paragraph{Results.} Table~\ref{tab:e1} gives the raw counts.
\PctUnderBound{} \CIUnderBound{} of the \NumStateChangingOps{} state-changing
operations in the corpus sit on a collided key. Nearly all of this comes from
AWS's query and RPC protocols: \NumMOneOpsAWS{} of the \NumMOneOps{} M1 flags
are AWS operations, and EC2 alone puts \NumEcTwoOpsPerKey{} operations on each of
\texttt{POST /} and \texttt{GET /}. AWS itself is not exposed, because SigV4
signs the query and \hdr{X-Amz-Target}; it shows how common the protocol pattern
is. Outside AWS, \PctApisCollidedNonAWS{} \CIApisCollidedNonAWS{} of APIs are
flagged, and the hand-confirmed figure is \NumConfirmedApisNonAWS{} of
\NumSpecsNonAWS{} (\PctConfirmedApisNonAWS); the true share lies between the
two, nearer the lower end given the M2 and M4 precision. A census of all
\NumFinanceFlagged{} finance flags confirmed \NumFinanceTP{}, in
\NumFinanceApisTP{} of \NumSpecsFinance{} finance APIs. The confirmed
collisions select high-impact commands: Stripe's usage records
(\texttt{action=increment|set}), Zuora refunds (\texttt{settle|reject}),
Zoom's \texttt{DELETE /users/\{id\}?action=disassociate|delete}, and openSUSE's \texttt{POST /person/\{login\}}, whose \texttt{cmd} parameter selects \texttt{change\_password}, \texttt{lock} or \texttt{delete}. Only
\NumApisDocOverride{} APIs (\PctApisDocOverride, all Amadeus) document a method
override, which is a framework property that descriptions rarely state
(\S\ref{sec:prev:e2}). Using every spec version instead of one per API changes
the headline share to \PctUnderBoundAllFiles{} (\PctUnderBoundAllFilesNonAWS{}
excluding AWS).

\paragraph{Standards that under-bind by construction.} Two specifications
outside the corpus make \GA the normal case. In HL7 FHIR, conditional
update and delete select their targets by query, and a \texttt{POST [base]}
transaction Bundle carries every CRUD interaction in its body; the public HAPI
FHIR R4 server declares conditional update, create and delete for all
\NumFhirTypesConditional{} resource types. The Model Context Protocol sends every
tool call as a POST to one endpoint and selects the operation by the JSON-RPC
body. Its 2026 revision mirrors the method and tool name into headers and
requires servers to reject mismatches, citing exactly the split-view risk of
\GB; DPoP does not occur in the MCP specification files we pinned. A proof
for either protocol names an endpoint, not an operation.

\subsection{Framework method override (\Exp{2})}\label{sec:prev:e2}

\paragraph{Method.} For \NumTargetFrameworks{} target frameworks (and Symfony,
Jersey, Hono and gorilla/handlers for comparison) we read the override code at a
pinned release tag and the vendor documentation. All \NumPinnedSourcesVerified{}
pinned source files were re-fetched and compared byte for byte, and all
\NumDocChecks{} documentation claims were checked against the fetched pages. To
estimate how often projects enable an opt-in mechanism, we use GitHub code
search counts; these are engine estimates over files, so we report them only as
orders of magnitude.

\paragraph{Results.} Override is reachable in the default configuration of
\NumOverrideFrameworks{} of the \NumTargetFrameworks{} frameworks. Rails keeps
\texttt{Rack::MethodOverride} unless an app is API-only, and Laravel enables
Symfony's override on every request, including the
\hdr{X-HTTP-METHOD-OVERRIDE} header for any POST. It is opt-in in
\NumOptInFrameworks{} (Spring Boot since 2.2, ASP.NET Core, Express through the
\texttt{method-override} package) and absent in \NumNoOverrideFrameworks{}
(FastAPI/Starlette, Django, Go \texttt{net/http} and chi, where the only method
change is HEAD to GET). Every built-in mechanism rewrites only POST, so the
transition is always \emph{a proof for POST executing PUT, PATCH or DELETE}.
The variant that reads \texttt{\_method} from the \emph{query} is the worst
case for DPoP, since \htu omits the query: it is reachable in Spring MVC,
Express with \texttt{methodOverride('\_method')}, gorilla/handlers,
Laravel/Symfony and Jersey. ASP.NET Core accepts any override value, with no
allowlist.

Where override is opt-in, roughly 1--2\% of projects turn it on (Spring
\PctEnableSpring{}, Express \PctEnableExpress{} of \texttt{package.json} files,
\PctEnableExpressNpm{} by npm downloads, gorilla \PctEnableGorilla{}); ASP.NET
Core is lower (\PctEnableAspnet{}). Where it is on by default it stays on: only
\PctRailsApiOnly{} of Rails application files declare API-only mode, so about
\PctRailsOverrideOn{} keep the middleware, and Laravel has no switch. Among
Express projects that use the package, \PctExpressQueryGetter{} of the two
common call forms choose the query getter.

\paragraph{Limitations.} The corpus over-represents hyperscalers and
under-represents GraphQL, gRPC and private APIs; recall is unmeasured, so
counts are lower bounds; one author judged the sample. A collision means a
proof bound to $K$ would be ambiguous, not that the API accepts DPoP.
Code-search ratios are engine estimates. Override is exploitable only if a
POST route shares its path with a state-changing route and the verifier sits
in front of the override, which the harness tests (\S\ref{sec:results}).
